\documentclass[11pt]{article}
\usepackage{mathunal-base}
\mathunalfuente{serif}
\newlist{opciones}{enumerate}{1}
\setlist[opciones]{label=(\alph*),itemsep=1pt,topsep=2pt,leftmargin=1.8em,font=\normalfont}

\renewcommand{\examtitulo}{Primer Parcial, Métodos Numéricos}
\renewcommand{\examfecha}{Marzo 10 de 2007}

\begin{document}

\examheader

\vspace{6pt}
\noindent Nombre \dotfill\ Carné \dotfill\ Grupo \dotfill

\vspace{8pt}
\noindent\textbf{La interpretación del examen hace parte de la evaluación, por tanto no se admiten preguntas.}

\vspace{10pt}
\noindent\textbf{1.} Considere el proceso de iteración de punto fijo $x_{n+1}=g(x_n)$ con $g(x)=\dfrac{1}{3}x-5$.
\begin{opciones}
\item (5\%) Pruebe que $\alpha=-\dfrac{15}{2}$ es un punto fijo de $g$.
\item (10\%) Pruebe que $|\alpha-x_{n+1}|=\dfrac{1}{3}|\alpha-x_n|$ para $n=0,1,\ldots$
\item (5\%) Pruebe que $|\alpha-x_{n+1}|=\dfrac{1}{3^{n+1}}|\alpha-x_0|$ para $n=0,1,\ldots$
\item (5\%) Con base en lo anterior, justifique la siguiente afirmación: Esta iteración de punto fijo converge para cualquier aproximación inicial $x_0$.
\end{opciones}

\vspace{5mm}
\noindent\textbf{2.} Consideremos el sistema no lineal
\[
\begin{aligned} x^2+y^2&=2\\ xy&=1 \end{aligned}
\]
\begin{opciones}
\item (5\%) Compruebe que las soluciones del sistema son: $\begin{bmatrix}1\\1\end{bmatrix}$ y $\begin{bmatrix}-1\\-1\end{bmatrix}$.
\item (10\%) Realice una iteración del método de Newton con primera aproximación $x^{(0)}=\begin{bmatrix}0\\2\end{bmatrix}$.
\item (10\%) Justifique el siguiente enunciado: Para este sistema, el método de Newton no se puede aplicar con cualquier aproximación inicial.
\end{opciones}

\vspace{5mm}
\noindent\textbf{3. (25\%)} Indique si los siguientes enunciados son verdaderos o falsos y justifique su respuesta.

\noindent Sean $A=\begin{bmatrix}2&3\\1&4\end{bmatrix}$ y $b=\begin{bmatrix}5\\5\end{bmatrix}$.
\begin{opciones}
\item \rule{2cm}{0.4pt} El sistema $Ax=b$ tiene una única solución.
\item \rule{2cm}{0.4pt} Si $x^{(0)}=\begin{bmatrix}0\\0\end{bmatrix}$, entonces la iteración de Gauss-Seidel para aproximar la solución de $Ax=b$ converge.
\item \rule{2cm}{0.4pt} $\|A\|_1=\|A\|_\infty$.
\item \rule{2cm}{0.4pt} Si las dos componentes de un vector $z$ son iguales, entonces $z$ es una solución de $Ax=b$.
\item \rule{2cm}{0.4pt} La iteración de Jacobi para aproximar la solución de $Ax=b$ no converge para cualquier vector inicial $x^{(0)}$.
\end{opciones}

\vspace{5mm}
\noindent\textbf{4.} Sea $f(x)=e^{2x}-e^x-2$.
\begin{opciones}
\item (10\%) Pruebe que $f$ tiene una única raíz en el intervalo $[0,1]$.
\item (5\%) Realice una iteración del método de Newton con primera aproximación $x_0=0$.
\item (10\%) Pruebe o refute la afirmación siguiente: Para encontrar esta raíz, el método de Newton converge cuadráticamente.
\end{opciones}

\vfill
\rule{\linewidth}{0.4pt}\\[1pt]
{\scriptsize\textit{Transcripción digital --- MathUNAL}\hfill\textcolor{unalblue}{\textbf{1}}}

\end{document}
